% ============================================================
% 绪论（教材：戈宝军《电机学》第三版，绪论，印刷页 1--2）
% 教材中绪论不编章号，故用 \chapter* 保持"第X章"与教材一致
% ============================================================
\chapter*{绪论}
\markboth{绪论}{}
\phantomsection\addcontentsline{toc}{chapter}{绪论}

电机学研究的对象是\textbf{电机}——以电磁感应定律为基础、实现机械能与电能之间转换以及电能特性变换的机械装置。它包含两大类：做相对运动的\textbf{旋转电机}与不动的\textbf{变压器}。本课程的任务不是记住一堆结构图，而是建立一套"磁场—磁路—电路—能量"的分析框架，并用它去解释四类电机为什么这样运行、性能由什么决定、参数如何测定。

这门课的特点，教材在前言里说得直白：理论性强、概念多、电磁与时空关系复杂、与工程实际联系紧密。因此在读每一章时都要反复问三个问题：\textbf{磁场从哪来}（励磁）、\textbf{磁场怎么动}（相对运动或交变）、\textbf{能量怎么过去}（电磁功率与损耗）。这三个问题一旦答不上来，后面的等值电路就只是符号游戏。

% ------------------------------------------------------------
\section*{一、电机的概念}
\addcontentsline{toc}{section}{一、电机的概念}

电机是以电磁感应定律为基础，实现机械能与电能之间的转换，以及电能特性变换的机械装置，包括\textbf{旋转电机}和\textbf{变压器}两大类。

\begin{itemize}
  \item \textbf{旋转电机}是机电能量转换装置：作为\textbf{发电机}运行时把机械能转换成电能；作为\textbf{电动机}运行时把电能转换成机械能；有的电机还用作\textbf{调相机}，用于改善电网的功率因数；此外还有大量\textbf{微特电机}，广泛用于自动控制系统中。
  \item \textbf{变压器}是各部件间无相对运动的电能变换装置，广泛用于电能传输，以及电压、电流、阻抗的变换和电路隔离。
\end{itemize}

\begin{center}
\begin{tabular}{@{}lll@{}}
\toprule
装置 & 有无相对运动 & 能量/信号的作用 \\
\midrule
发电机 & 有 & 机械能 $\rightarrow$ 电能 \\
电动机 & 有 & 电能 $\rightarrow$ 机械能 \\
调相机 & 有 & 只交换无功，改善功率因数 \\
变压器 & 无 & 电压、电流、阻抗变换，电路隔离 \\
控制电机 & 有/无 & 控制系统中作执行、检测、放大、解算元件 \\
\bottomrule
\end{tabular}
\end{center}

\corepoint{发电机与电动机只是同一台电机的两种\textbf{运行方式}，电机本身是\textbf{可逆}的。判断某台电机处于发电机态还是电动机态，看的不是它的构造，而是电磁转矩与转速的方向关系、以及功率的流向（第\,12、16、20\,章会反复用到这条判据）。}

% ------------------------------------------------------------
\section*{二、电机的主要作用}
\addcontentsline{toc}{section}{二、电机的主要作用}

教材把电机的作用归为三个方面。

\subsection*{1. 电能的生产、传输和分配}
在发电厂，通过汽轮机、燃气轮机、柴油机、水轮机、风（轮）机等的驱动，把燃料燃烧和原子核裂变产生的能量，或水的势能、风的动能等转化为机械能，由\textbf{发电机}完成机电能量转换生产电能；然后经\textbf{变压器}升高电压，通过输电线路送到各用电地区，再经变压器降低电压供用户使用。

这一条链上，电机学要回答的是：为什么输电要升压（同样的功率下，电压升高则电流减小、线路损耗按电流平方下降）；为什么发电机出口电压受限于绝缘与制造（第\,16 章）；为什么变压器容量与铁心截面、绕组匝数直接相关（第\,3--4 章）。

\subsection*{2. 驱动各种生产机械和装备}
在工农业、交通运输、国防等部门和日常生活中，电动机广泛用于驱动生产机械、设备和器具：机床驱动、电力给排水、农副产品加工、矿石采掘和传送、电驱动车的牵引、鼓风设备、起重设备、轧钢机械、造纸设备、化工机械及家用电器的驱动等。全社会的用电量中，电动机（尤其是感应电动机）占绝大部分，这也是本课程把感应电机放在核心位置的原因。

\subsection*{3. 自动控制系统中的控制元件}
控制电机在控制系统、自动化系统中作为执行、检测、放大和解算元件。这类电机功率一般较小，但品种繁多、用途广泛（步进电机、伺服电机、测速发电机、旋转变压器等）。

% ------------------------------------------------------------
\section*{三、电机的分类}
\addcontentsline{toc}{section}{三、电机的分类}

电机种类很多，结构和用途各不相同，分类方法也很多。电机学中常采用以下三种分类方法。

\subsection*{（1）按应用的电流种类}
分为\textbf{直流电机}和\textbf{交流电机}。

\subsection*{（2）按能量转换功能和用途}
\begin{enumerate}
  \item 将机械能转换成电能的\textbf{发电机}；
  \item 将电能转换成机械能的\textbf{电动机}；
  \item 将一种交流电压等级改变成另一种交流电压等级的\textbf{变压器}；
  \item 在控制系统中起执行和测量作用的\textbf{控制电机}。
\end{enumerate}
另外还有\textbf{变流机}、\textbf{变频机}和\textbf{移相机}，分别用来改变电流、频率和相位。

\subsection*{（3）按运行转速的特点}
\begin{enumerate}
  \item \textbf{静止设备}的变压器；
  \item 没有固定同步转速的\textbf{直流电机}；
  \item 转子转速与同步转速有差异的\textbf{感应电机}；
  \item 转速恒为同步转速的\textbf{同步电机}。
\end{enumerate}

随着电力电子技术和电工材料的发展，出现了不属于上述传统类型的电机，如步进电机、无刷电机、超声波电机等，通常称为\textbf{特种电机}。本书按\textbf{工作原理}介绍电机，主要包括变压器、感应电机、同步电机和直流电机四大类电机的基本原理与应用。

\begin{center}
\begin{tabular}{@{}p{2.4cm}p{2.2cm}p{2.6cm}p{5.0cm}@{}}
\toprule
大类 & 电流种类 & 转速特点 & 本课程关注的核心问题 \\
\midrule
变压器 & 交流 & 静止 & 变比、等值电路、参数测定、联结组、并联运行 \\
感应电机 & 交流 & $n \neq n_1$（异步） & 转差率、等值电路、转矩—转差特性、起动与调速 \\
同步电机 & 交流 & $n = n_1$（同步） & 电枢反应、功角特性、并网、功率调节 \\
直流电机 & 直流 & 无固定同步转速 & 电枢反应、感应电动势与电磁转矩、机械特性 \\
\bottomrule
\end{tabular}
\end{center}

\corepoint{分类的目的不是背名词，而是确定"用哪一套方程"。看到"转速恒等于同步转速"就要想到同步电机方程（$E_0$ 与 $U$ 的关系、功角 $\delta$）；看到"转速随负载变化"就要想到感应电机的转差率 $s$；看到"直流"就要想到换向器把交流整成直流、以及 $E_a=C_e\Phi n$ 这条线性关系。}

% ------------------------------------------------------------
\section*{四、电机的发展}
\addcontentsline{toc}{section}{四、电机的发展}

英国物理学家法拉第（M. Faraday）于 1821 年发现了载流导体在磁场中会受到力的作用的现象，1831 年又发现了电磁感应定律，并很快就出现了原始模型电机。从此，电机的研究和应用迅速发展起来，至今已近 200 年。电机发展的历史大体上可以分成四个时期：

\begin{enumerate}
  \item \textbf{直流电机的产生和发展时期}；
  \item \textbf{交流电机的发展时期}；
  \item \textbf{电机理论、设计和制造工艺逐步达到完善化的时期}；
  \item \textbf{新技术、新材料等与电机相结合发展的新时期}。
\end{enumerate}

\begin{center}
\begin{tabular}{@{}llp{8.2cm}@{}}
\toprule
年份 & 人物/事件 & 意义 \\
\midrule
1821 & 法拉第发现载流导体受力 & 电动机原理的起点 \\
1831 & 法拉第发现电磁感应定律 & 发电机与变压器的共同理论基础 \\
1869 & 格拉姆（Z. T. Gramme）制成实用直流发电机 & 电力发展的重要推动 \\
1882 & 爱迪生（T. A. Edison）建成纽约曼哈顿直流照明系统 & 直流供电进入商业运行 \\
1883 & 特斯拉（N. Tesla）研制出第一台两相感应电机 & 交流电机进入实用阶段 \\
1888 & 多利沃-多勃罗夫斯基制成三相感应电机 & 三相制的基础 \\
20 世纪初 & 三相交流制在电力工业中占据统治地位 & 同步电机分高速（汽轮）与低速（水轮）两类 \\
\bottomrule
\end{tabular}
\end{center}

电机发展初期主要是直流电机发展的历史。随着直流电机的广泛应用，它很快暴露出固有缺点：不能解决远距离输电以及电压高低变换的问题；随着单机容量日益增大，换向也越来越困难，于是交流电机获得了迅速发展。19 世纪 90 年代初开始，三相同步电机按结构逐渐划分为高速和低速两类，高速的以汽轮发电机为代表，低速的以水轮发电机为代表。到 19 世纪末，直流电机、感应电机、同步电机、变压器等常规电机都已得到迅速发展和应用，相应的基本理论和设计方法也已初步建立。

20 世纪是电机发展的新时期。就单机容量而言，水轮发电机的最大单机容量在 20 世纪初不超过 1\,MW，现在已达到 1000\,MW；汽轮发电机的单机容量在 20 世纪初不超过 5\,MW，1930 年提高到 100\,MW，20 世纪 40—50 年代由于采用了氢冷、氢内冷、油冷和水冷等冷却方法，单机容量进一步提高，目前已超过 1000\,MW。由于新的永磁材料的出现，制成了许多高效节能、维护简单的永磁电机；由于电机和电力电子装置、单片机相组合，出现了各种性能和形态迥异的"一体化电机"。交流调速电机、风力发电机、核电电机和高效节能电机等方面的发展最为令人瞩目。

% ------------------------------------------------------------
\section*{五、本课程的任务}
\addcontentsline{toc}{section}{五、本课程的任务}

本课程是一门专业基础课。通过学习可获得电机的基本理论、基本知识和基本技能，为学习专业课做好准备，也为今后从事有关的专业工作打下基础。学习本课程后应达到下列基本要求。

\begin{enumerate}
  \item 对\textbf{磁路}的基本概念、基本定律、计算方法及交流铁芯线圈的励磁特性有基本的了解。
  \item 对变压器和三种主要电机（感应电机、同步电机和直流电机）的\textbf{基本结构}，交流三相单层、双层绕组和直流单叠、单波绕组的连接规律，要有一定认识；对各种电机中定、转子的\textbf{磁动势和气隙磁场}的性质及时间、空间关系，要有深入了解。
  \item 对变压器和三种主要电机正常\textbf{稳态运行}时的分析方法和运行性能，要牢固掌握：能正确地建立基本方程，弄清电机中的能量传递和能量转换关系；对稳态运行时的参数要有清晰的物理概念，能熟练运用\textbf{等值电路}计算变压器和交流电机的性能。
  \item 了解电机的有关\textbf{工程性问题}，如发热、冷却、温升、励磁系统和各种电机的应用范围；掌握电机的额定值、变压器和电机的主要性能数据及范围。
  \item 了解常用的特殊变压器、特殊感应电机和同步永磁电动机的基本结构、特点与应用。
  \item 通过\textbf{实验}，熟练掌握电机的基本实验方法和操作技能。
\end{enumerate}

\subsection*{学习方法的四条建议}
\begin{enumerate}
  \item \textbf{先建方程，再画图。}四类电机的分析套路几乎相同：写电压方程 $\rightarrow$ 写磁动势（磁通）方程 $\rightarrow$ 归算 $\rightarrow$ 画等值电路 $\rightarrow$ 画相量图 $\rightarrow$ 由等值电路求性能。顺序颠倒就会变成"背图"。
  \item \textbf{抓住"磁动势平衡"这一条主线。}变压器是一次、二次电流产生的磁动势共同建立主磁通；感应电机是定子、转子磁动势合成建立气隙磁场；同步电机是励磁磁动势与电枢磁动势合成。三者形式不同，实质相同（第\,4、13、17\,章）。
  \item \textbf{把"归算"当成一种坐标变换。}归算不改变功率和损耗，只改变电压、电流、阻抗的数值，目的是把两个不同电压等级（或不同频率的转子）接到同一个电路里（第\,4.3 节、第\,13.3 节）。
  \item \textbf{每个结论都要能说出前提。}如"$T \propto U^2$"只在忽略定子阻抗、转差率不变的近似下成立；"$U_1\approx E_1$"只在忽略一次漏阻抗时成立。写公式时把条件一并写上，考试与工程都不会吃亏。
\end{enumerate}

\corepoint{本课程的计算题几乎全部落在三件事上：\textbf{参数}（由空载、短路试验求 $R_m$、$X_m$、$R_k$、$X_k$ 等）、\textbf{特性}（由等值电路求 $I$、$P$、$T$、$\eta$、$\Delta U$）、\textbf{运行方式}（起动、调速、制动、并联、并网）。复习时按这三条把每章内容重新归类一遍，比按页码顺序重读有效得多。}
